\documentclass[10pt,a4paper]{amsart}
\usepackage[margin=19mm]{geometry}
\usepackage{amsmath,amssymb,mathtools}
\usepackage[scr=rsfs,cal=euler]{mathalfa}
\usepackage{xfrac}
\usepackage[unicode,hidelinks,bookmarks=false]{hyperref}
\newcommand{\ie}{i.e.\ }
\newcommand{\cf}{cf.\ }
\newcommand{\resp}{respectively\ }
\newcommand{\Subsection}{\subsection}
\providecommand{\nobreakdash}{\nobreak\hbox{-}}
\setlength{\emergencystretch}{2em}
\multlinegap0pt
\usepackage{bm}
\usepackage{bbm}
\usepackage{dsfont}
\usepackage{xspace}
\usepackage{cancel}
\usepackage{mathtools}
\usepackage{amsthm}
\makeatletter
\@ifundefined{theorem}{\newtheorem{theorem}{Theorem}}{}
\@ifundefined{proposition}{\newtheorem{proposition}[theorem]{Proposition}}{}
\makeatother
\title[Factorization of Isomorphisms of $(H,\theta)$-twisted Lie al]{Factorization of Isomorphisms of $(H,\theta)$-twisted Lie algebroids}
\author{Nasser Saipele nansidi, Bertuel Tangue Ndawa, Luc Emery Diekouam Fotso, Emmanuel Fouotsa, Joseph Dongho}
\date{Source: arXiv:2608.09306v2 (CC BY 4.0)}
\begin{document}
\maketitle

\noindent\emph{Source and licence.} Factorization of Isomorphisms of $(H,\theta)$-twisted Lie algebroids. Version arXiv:2608.09306v2 (CC BY 4.0), distributed under the Creative Commons Attribution 4.0 International licence, as recorded in the arXiv licence metadata for this exact version and shown on the abstract page https://arxiv.org/abs/2608.09306v2. Full rights notice: licenses/arxiv-2608.09306-CC-BY-4.0.txt.\par\smallskip

\noindent\textit{Editorial scope.} Counted unit: the Proposition whose source label is \texttt{prop:Courant-theta} in the article (source0.tex lines 325--335, source label \texttt{prop:Courant-theta}), delivered together with the article's own complete original proof of it (source0.tex lines 337--402, one \texttt{proof} environment of 66 lines). No mathematics is written, repaired or re-derived: statement and proof are the article's own text line for line. Required context retained uncounted: cartan1 (lines 305--309); cartan2 (lines 305--309); cartan3 (lines 305--309). Every definition, display and label of those lines is retained unchanged. Omitted from this package and not redistributed: 1--304, 310--324, 336--336, 403--1290. Nothing inside a retained excerpt is omitted or altered: every retained body is delivered exactly as the article's lines are written, so no omission receipt of any kind accompanies this package. Citations: the retained excerpts carry no citation command at all, so no bibliography entry of the article is transcribed and no reference is redistributed. Reference adaptation: none was needed and none was made; every reference inside the delivered unit resolves inside it, so no unresolved reference or citation is produced. Printed slips kept literal: no printed word, symbol, hypothesis or number of the counted statement or of its proof was altered. Layout and class: the article's own class is \texttt{article} with packages including xcolor, amsmath, amsthm, amssymb, amsfonts, mathtools, geometry, hyperref, enumitem, bm, fontenc, authblk, indentfirst, comment; the wrapper is the lane's pinned \texttt{amsart} editorial wrapper at 10pt on A4, which restates the theorem-like environments the retained text needs and adds the article's own macro definitions that the retained text needs (none), with extra wrapper packages (bm, bbm, dsfont, xspace, cancel, mathtools). The delivered numbering is the unit's own. Assets: no retained span contains a figure environment, a graphics-inclusion command, a PDF-page inclusion, a table or a TikZ picture; no image file is copied into this package, the package declares no asset dependency and no image file is redistributed with it. Licence: version arXiv:2608.09306v2 of this article is distributed under the Creative Commons Attribution 4.0 International licence, as recorded in the arXiv licence metadata for this exact version and shown on the abstract page https://arxiv.org/abs/2608.09306v2. A full rights notice is delivered at licenses/arxiv-2608.09306-CC-BY-4.0.txt. Characters: every retained excerpt is pure ASCII, so the delivered engine, XeLaTeX with its default Latin Modern text font, reports no missing character.
\begin{align}
	[\mathcal{L}^{\theta}_{X},d_{\theta}] &= 0, \label{cartan1}\\
	[\mathcal{L}^{\theta}_{X}, i_{Y}] &= i_{[X,Y]}, \label{cartan2}\\
	[\mathcal{L}^{\theta}_{X},\mathcal{L}^{\theta}_{Y}] &= \mathcal{L}^{\theta}_{[X,Y]}. \label{cartan3}
\end{align}

\begin{proposition}\label{prop:Courant-theta}
 For all $e_i\in\Gamma(E)$ and $f\in C^\infty(M)$, the quadruple $(E,\rho,\langle\cdot,\cdot\rangle,\circ_\theta,D_\theta)$ satisfies de the following properties
\begin{enumerate}
\item $\rho(e_1\circ_\theta e_2) = [\rho(e_1),\rho(e_2)]$;
\item $e_1\circ_\theta e_2 + e_2\circ_\theta e_1 = D_\theta\langle e_1,e_2\rangle$;
\item $\langle e_1\circ_\theta e_2, e_3\rangle + \langle e_2, e_1\circ_\theta e_3\rangle = \mathcal{L}^\theta_{\rho(e_1)} \langle e_2,e_3\rangle$;
\item $e_1\circ_\theta (e_2\circ_\theta e_3) = (e_1\circ_\theta e_2)\circ_\theta e_3 + e_2\circ_\theta (e_1\circ_\theta e_3)$;
\item $D_\theta f \circ_\theta e_1 = 0$,\quad $e_1 \circ_\theta D_\theta f = D_\theta\langle e_1, D_\theta f\rangle$.
\end{enumerate}
We will term by \emph{$\theta$-twisted Courant algebroid } the quadruple $(E,\rho,\langle\cdot,\cdot\rangle,\circ_\theta,D_\theta)$ satisfying such properties.
\end{proposition}

\begin{proof}
\textit{Axiom 1:} The vector component of $e_1\circ_\theta e_2$ is $[X_1,X_2]$, which is exactly equal to $[\rho(e_1),\rho(e_2)]$.

\smallskip
\textit{Axiom 2:} Using $\mathcal{L}^\theta_X = d_{\theta} i_X + i_X d_{\theta}$, the $1$-form part of
$e_1\circ_\theta e_2 + e_2\circ_\theta e_1$ equals
\begin{align*}
	\bigl(d_{\theta} i_{X_1}\xi_2 + i_{X_1}d_{\theta}\xi_2 - i_{X_2}d_{\theta}\xi_1\bigr)
	&+ \bigl(d_{\theta} i_{X_2}\xi_1 + i_{X_2}d_{\theta}\xi_1 - i_{X_1}d_{\theta}\xi_2\bigr)\\
	&= d_{\theta}\bigl(i_{X_1}\xi_2 + i_{X_2}\xi_1\bigr)\\
	&= d_{\theta}\langle e_1,e_2\rangle.
\end{align*}
The vector parts cancel since $[X_2,X_1]=-[X_1,X_2]$, hence
$e_1\circ_\theta e_2 + e_2\circ_\theta e_1 = D_\theta\langle e_1,e_2\rangle$.

\smallskip
\textit{Axiom 3:} Expanding the left-hand side and using $i_{X_3}i_{X_2}=-i_{X_2}i_{X_3}$, one obtains
$$
i_{[X_1,X_2]}\xi_3 + i_{X_3}\mathcal{L}^\theta_{X_1}\xi_2
+ i_{X_2}\mathcal{L}^\theta_{X_1}\xi_3 + i_{[X_1,X_3]}\xi_2.
$$
On the other hand, applying $\mathcal{L}^\theta_{X_1}$ to
$\langle e_2,e_3\rangle = i_{X_2}\xi_3 + i_{X_3}\xi_2$, and using \eqref{cartan2},
we obtain exactly the same expression.

\smallskip
\textit{Axiom 4:} Set $\xi_{23}:=\mathcal{L}^\theta_{X_2}\xi_3 - i_{X_3}d_{\theta}\xi_2$, so that
$e_2\circ_\theta e_3 = \bigl([X_2,X_3],\ \xi_{23}\bigr)$.
Then
\begin{equation*}
%	\label{eq:lhsCA4}
	e_1\circ_\theta (e_2\circ_\theta e_3)
	=
	\Bigl([X_1,[X_2,X_3]],\ \mathcal{L}^\theta_{X_1}\mathcal{L}^\theta_{X_2}\xi_3
	- \mathcal{L}^\theta_{X_1}i_{X_3}d_{\theta}\xi_2
	- i_{[X_2,X_3]}d_{\theta}\xi_1\Bigr).
\end{equation*}
For the right-hand side, one computes $(e_1\circ_\theta e_2)\circ_\theta e_3$ and
$e_2\circ_\theta (e_1\circ_\theta e_3)$ separately.
The vector parts give
$[[X_1,X_2],X_3] + [X_2,[X_1,X_3]] = [X_1,[X_2,X_3]]$ by the Jacobi identity.
For the $1$-form parts, identity \eqref{cartan3} implies
$\mathcal{L}^\theta_{X_1}\mathcal{L}^\theta_{X_2}
= \mathcal{L}^\theta_{X_2}\mathcal{L}^\theta_{X_1}+\mathcal{L}^\theta_{[X_1,X_2]}$,
which matches the $\xi_3$-terms.
Applying \eqref{cartan2} to rewrite
$\mathcal{L}^\theta_{X_1}i_{X_3}
= i_{X_3}\mathcal{L}^\theta_{X_1}+i_{[X_1,X_3]}$,
and using the commutation $d_{\theta}\mathcal{L}^\theta_{X_1}=\mathcal{L}^\theta_{X_1}d_{\theta}$,
which follows from \eqref{cartan1},
all remaining terms cancel identically.
Hence Axiom 4 holds.

\smallskip
\textit{Axiom 5:} Since $D_\theta f=(0,d_{\theta} f)$ has vanishing vector part, we have
$D_\theta f\circ_\theta e=0$.
For the second identity, we compute
$$
e\circ_\theta D_\theta f
= \mathcal{L}^\theta_X d_{\theta} f
= d_{\theta}\bigl(i_X d_{\theta} f\bigr)
= D_\theta\bigl(i_X d_{\theta} f\bigr)
= D_\theta\langle e,D_\theta f\rangle.
$$
This completes the proof.
\end{proof}

\end{document}
